\documentclass[12pt]{article}
\usepackage[margin=1in]{geometry}
\usepackage{amsmath, amssymb, amsthm, graphicx, hyperref}
\usepackage{subfigure}
\usepackage{color}
\usepackage{fancyhdr}
\usepackage{multirow, multicol}
\usepackage{pdfpages}
\usepackage{stfloats}
\usepackage{tikz}
\pagestyle{fancy}

\fancypagestyle{firststyle}
{
  \renewcommand{\headrulewidth}{0pt}
  \fancyhead[LO]{\includegraphics{nyu_short_black}}
  \fancyhead[RO]{}
	\fancyfoot{}
}

\fancyhead[LO]{\scalebox{0.25}{\includegraphics{nyu_short_black}}}
\fancyhead[RO]{\tiny MATH-UA.120.002}


%\fancyfoot[RO]{\begin{tikzpicture}
%[scale=1, line/.style={black}]
%\draw[line] (0, 0)--(2,0)--(2, 2)--(0, 2)--(0, 0);
%%\node at (-0.5, 1) [anchor=east]{\large Score};
%\end{tikzpicture}}

\newcommand*\mycirc[1]{%
  \begin{tikzpicture}[baseline=(C.base)]
    \node[draw,circle,inner sep=1pt](C) {#1};
  \end{tikzpicture}}

\newcommand{\R}{\mathbb{R}}
\newcommand{\E}{\mathbb{E}}
\newcommand{\proj}{\textbf{proj}}
\newcommand{\curl}{\textrm{curl }}
\newcommand{\divergence}{\textrm{div }}
% For boldface vectors:
\renewcommand{\vec}[1]{\mathbf{#1}}
\renewcommand\arraystretch{2}

\newtheorem*{proposition}{Proposition}


\begin{document}
\thispagestyle{firststyle}
~
\vspace{0.5cm}

\begin{center}
\textbf{\Large MATH-UA.120.002 - Discrete Mathematics \\ \vspace{0.5cm} Midterm 2}\\~\\
Thursday, March 31, 2016
\end{center}

\vspace{1cm}

%\noindent
%\textbf{Name:} \underline{\hspace{15cm}}

\vspace{1cm}

\noindent
This exam is scheduled for 110 minutes, to be done individually, without calculators, notes, textbooks, and other outside materials.
\vspace{0.5cm}

\noindent
\textbf{Show all work to receive full credit, except where specified.}

\vspace{1cm}

\begin{multicols}{2}
\noindent

\noindent
\fbox{\parbox{10cm}{
\vspace{1cm}
\noindent
\textbf{Name: }

\vspace{1cm}

\noindent
\textbf{NYU NetID (email): }


\vspace{2cm}

\noindent
\emph{I pledge that I have observed the NYU honor code, and that I have neither given nor received unauthorized assistance during this examination.}
\noindent



\noindent

\vspace{1cm}

\noindent
\textbf{Signature:}

\vspace{2cm}
}}
%
%\vspace{0.5cm}
%

\begin{flushright}
\begin{tabular}{| c | c |}
\hline
Problem & Points \\
\hline
% & \\
TF & ~~~~~~~~~/20\\[0.5cm]
%& \\
\hline
% & \\
1 & ~~~~~~~~~/20\\[0.5cm]
%& \\
\hline
% & \\
2 & ~~~~~~~~~/20\\[0.5cm]
%& \\
\hline
% & \\
3 & ~~~~~~~~~/20\\[0.5cm]
%& \\
\hline
% & \\
4 & ~~~~~~~~~/20\\[0.5cm]
%& \\
\hline
\hline %& \\
Total & ~~~~~~~~~/100\\[0.5cm]
%& \\
\hline
\end{tabular}
\end{flushright}

\end{multicols}

\newpage

\section*{Answer Sheet (for Problem 1 True/False)} \label{answersheet}

%(3 points each) Select the correct option.  You need not explain your answer.  No partial credit will be given.
%
%\vspace{0.5cm}

\noindent
\textbf{You must record your final answers on the following answer sheet.  Only this page will be graded.}

\vspace{1cm}

\begin{center}
\begin{tabular}{| c | c c c c c |}
\hline

1 & \mycirc{T} & \mycirc{F} & & &  \\
% & & & & & \\
2 & \mycirc{T} & \mycirc{F} & & &   \\
% & & & & & \\
3 & \mycirc{T} & \mycirc{F} & & &   \\
% & & & & & \\
4 & \mycirc{T} & \mycirc{F} & & &  \\
% & & & & & \\
5 & \mycirc{T} & \mycirc{F} & & &  \\
6 & \mycirc{T} & \mycirc{F} & & & \\
7 & \mycirc{T} & \mycirc{F} & & & \\
8 & \mycirc{T} & \mycirc{F} & & & \\
9 & \mycirc{T} & \mycirc{F} & & & \\
10 & \mycirc{T} & \mycirc{F} & & & \\
\hline
\end{tabular}
\end{center}
\begin{flushright}
\begin{tikzpicture}
[scale=1, line/.style={black}]
\draw[line] (0, 0)--(2,0)--(2, 2)--(0, 2)--(0, 0);
\node at (-0.5, 1) [anchor=east]{\Large TOTAL};
\end{tikzpicture}
\end{flushright}
\renewcommand\arraystretch{1}

\begin{center}
\textbf{Make sure that you have marked your final answers on the answer sheet on page \pageref{answersheet}. Only answers on that page will be graded.}
\end{center}

\newpage
\section*{True or False}

(2 points each) Determine if each of the following statement is TRUE or FALSE.\\Choose TRUE if the statement is always true, and choose FALSE if it is not always true. Each question is worth 3 points.

\vspace{0.2cm}

\noindent
\textbf{Indicate your solution in the answer sheet on page \pageref{answersheet}.  You need not provide any justification.}


\begin{enumerate}
\item Ten points are placed in the plane with no two on the same horizontal line and no two on the same vertical line. Prove that there must be four points that can be joined to give an upward path or four points that can be joined to give a downwards path. (see figure)

OLDS:
\item The set $(0,1]\cup\mathbb{Z}$ satisfies the well ordering principle.
\vspace{0.8cm}
\item Let $A=\{1,2,3\}$ and $B=\{3,4,5,6\}$. Then the relation $\{(1,3), (2,5), (3,4)\}$ is a function $f:A\to B$.
\vspace{0.8cm}
\item Let $A$ and $B$ be finite sets. Assume $f:A\to B$ is one-to-one and that there is a function $g:B\to A$ that is one to one. Then $f$ must also be onto.
\vspace{0.8cm}
\item $(h\circ g\circ f)^{-1}=h^{-1}\circ g^{-1}\circ f^{-1}$
\vspace{0.8cm}
\item $(a-b)(a-c)(b-c)$ is always even if $a,b,c$ are integers.
\vspace{0.8cm}
\item Let $f(n)=n^2-3n+2$. Then $f(n)$ is $\Omega(n^3)$.
\vspace{0.8cm}
\item Let $A$ be a set. Suppose $f$ and $g$ are functions $f:A\to A$ and $g:A\to A$ with the property $f\circ g=\textrm{id}_A$. Then $f=g^{-1}$. (Warning: $A$ is not necessarily a finite set.)
\vspace{0.8cm}
\item Suppose $S$ is a set of $n + 1$ integers. There exist distinct $a, b \in S$ such that $a - b$ is a multiple of $n$.
\vspace{0.8cm}
\item Suppose $f(n)$ is $O(g(n))$. Then $34 f(n)$ is also $O(g(n))$.
\vspace{0.8cm}
\item Suppose $f,g:\mathbb{R}\to\mathbb{R}$ are defined by
\[
f(x)=x^2-4x+2,\qquad g(x)=2x-1
\]
Then $(f\circ g)(x)-(g\circ f)(x)=2x^2-4x+4$.
\end{enumerate}
\begin{center}
\textbf{Make sure that you have marked your final answers on the answer sheet on page \pageref{answersheet}. Only answers on that page will be graded.}
\end{center}
\newpage
\section*{Free Response}

\textbf{Show all work and justification.} 

\begin{enumerate}
\item (20 points) Solve the following recurrence relations.
\begin{enumerate}
\item $a_n=2-4a_{n-1}$; $a_0=2$.
\item $a_n=10a_{n-1}-25a_{n-2}$; $a_0=1$, $a_1=-3$.
\end{enumerate}
\newpage
\item (20 points) Suppose $f : A\to B$ and $g : B\to C$ are functions and $g\circ f$  is a one-to-one function.
Must f be one-to-one? Must g be one-to-one? Support your answers.
\newpage
\item (20 points)
\begin{enumerate}
\item (10 points) Use the method of smallest counterexample to show that for all $n\in\mathbb{N}$,
\[
\sum_{i=1}^n\frac{2}{(i+1)(i+2)}=\frac{n}{n+2}
\]
\item (10 points) Use the principle of induction to show that for all $n\geq 1$
\[
 1 \cdot 2 + 2 \cdot 3 + . . . + (n - 1) \cdot n =\frac{n-1}{n}
\]
WRONG PROBLEM
\end{enumerate}
\newpage
\item For a function $f : A\to B$, consider the induced inverse-image function $f^{-1}:2^B\to 2^A$ given by 
\[
f^{-1}(X)=\{a\in A: f(a)\in X\},\qquad X\in 2^B.
\]
\begin{enumerate}
\item (5 points) Let $f: \mathbb{N}\to \mathbb{N}$ be defined by $f(n)$ being the number of letters in the English word form of n. For example $f(1)=3$ (three letters in `one'), $f(4)=4$ (four letters in `four'). Find $f^{-1}(\{1,2,3,4\})$.

\item (15 points) Prove the follwing proposition. Note that this is NOT about the particular example in part (a).

\begin{proposition}
If $f:A\to B$ is a function and $f^{-1}:2^B\to 2^A$, defined as above, is onto, then $f$ is one-to-one. 
\end{proposition}
Hint: Prove the contrapositive.
\end{enumerate}

\end{enumerate}

\newpage
\noindent
\textbf{Extra paper}
%
%\newpage
%\noindent
%\textbf{Extra paper}

%\includepdf[pages=1-2]{FormulaSheets.pdf}

\end{document}